\section{Local definitions and choices}

\subsection{Memory State, cut, and present identity}
\begin{definition}[Memory State, inherited]
\label{def:memory-state}
The inherited Memory State is the AFC object
\[
p:V\to[0,1], \qquad \sum_{i\in V} p(i)=1.
\]
\end{definition}

\begin{definition}[Cut]
\label{def:active-cut}
For a declared Region $R\subseteq V$ with boundary $\partial R$, a cut is the declared inside/boundary/outside split $(R,\partial R,V\setminus R)$ on which the imported AFC identity is read.
\end{definition}

\begin{definition}[Present identity]
\label{def:present-identity}
The present identity is the AFC Stokes identity evaluated on the boundary of the declared cut.
\end{definition}

\subsection{\texorpdfstring{$\tau$}{tau} readout, frame/window, and active/dwell split}
\begin{choice}[$\tau$ readout]
\label{choice:tau-readout}
A declared readout map assigns a ternary symbol $\tau_t\in\{-1,0,+1\}$ to each update on the declared cut. The ternary readout is attached to that cut through the declared frames and windows while the declared cut remains fixed.
\end{choice}

\begin{definition}[Frame]
\label{def:frame}
A frame $f_t$ is the one-\(\mathrm{bip}_0\) readout slice on a declared cut unless a different working cut level is declared.
\end{definition}

\begin{definition}[Window]
\label{def:window}
A window $\omega$ is a finite ordered family of frames on a fixed declared readout contract.
\end{definition}

\begin{definition}[Active indicator and active/dwell split]
\label{def:active-dwell}
The active indicator is
\[
\chi_t:=\mathbf{1}\{\tau_t\in\{-1,+1\}\}.
\]
Active length counts the ticks with $\chi_t=1$ and dwell length counts the complementary ticks with $\tau_t=0$ on the same readout.
\end{definition}

\begin{definition}[Step flux and window current]
\label{def:step-window-current}
Step flux is the one-frame boundary-crossing quantity on a declared frame. Window current is the derived window-level accumulation of step flux over the ordered family of frames in $\omega$. When the retained transport layer is used, the windowed current-flux rate is namespaced as $\mathcal J_K$ and is kept distinct from the chirality-count rows $J_\alpha,J_\Omega,J_0$ and from the transverse witness $w_\perp$.
\end{definition}

\subsection{\texorpdfstring{$\tau$}{tau}-side readout, admissibility, and loss channel}
\begin{definition}[$\tau$ readout row]
\label{def:tau-readout-record}
The declared $\tau$-row records, on the declared frame/window family, the ternary symbols $\tau_t\in\{-1,0,+1\}$, the active/dwell partition, and any declared symbol counts $(N_+,N_0,N_-)$. These are row entries.
\end{definition}

\begin{definition}[Admissibility count and active set]
\label{def:mt}
At tick $t$, the declared gating rule produces an admissible successor set $A_t$. The admissibility count is
\[
m_t:=|A_t|\in\{1,2,3\}.
\]
When $m_t=1$, the realized update is fully constrained on that tick.
\end{definition}

\begin{definition}[Loss channel]
\label{def:loss-channel}
The admissibility-loss support on a declared window is the integer row
\[
L_{\mathrm{loss}}^{\mathbb Z}(\omega):=(n_1,n_2,n_3,T_{\mathrm{eval}}),
\qquad
n_m:=|\{t\in\omega\cap\mathrm{act}:m_t=m\}|,
\]
with
\[
T_{\mathrm{eval}}\in\mathbb Z_{>0},
\qquad
n_1+n_2+n_3\le T_{\mathrm{eval}}.
\]
The per-tick loss burden, accumulated burden, and declared window loss readout are
\[
u_t^{(\mathrm{loss})}:=\log_3\!\left(\frac{3}{m_t}\right),
\qquad
U^{(\mathrm{loss})}(\omega):=\sum_{t\in\omega\cap\mathrm{act}}u_t^{(\mathrm{loss})},
\]
\[
R_{\mathrm{loss}}(\omega):=u^{(\mathrm{loss})}(\omega):=\frac{1}{T_{\mathrm{eval}}}U^{(\mathrm{loss})}(\omega)
\quad[\mathrm{trit}/\mathrm{bip}_0].
\]
The equivalent multiplicative loss witness is
\[
\Gamma_{\mathrm{loss,raw}}(\omega):=\prod_{t\in\omega\cap\mathrm{act}}\frac{3}{m_t},
\qquad
\Gamma_{\mathrm{loss}}(\omega):=\Gamma_{\mathrm{loss,raw}}(\omega)^{1/T_{\mathrm{eval}}},
\qquad
T_\Gamma(\omega):=\Gamma_{\mathrm{loss}}(\omega)-1.
\]
Together, \(L_{\mathrm{loss}}^{\mathbb Z}\) and \(R_{\mathrm{loss}}\) name the loss channel on the declared window: the integer row records the admissibility support, and \(R_{\mathrm{loss}}\) is the retained trit-per-\(\mathrm{bip}_0\) loss readout assembled by the declared loss map.
\end{definition}

\begin{definition}[Active-only versus total duration count]
\label{def:rate-duration-count}
A rate row uses total ticks on the declared window unless the row declares an active-only duration count and records $T_{\mathrm{act}}$ alongside $T_{\mathrm{eval}}$.
\end{definition}

\subsection{Predecessor trace, successor horizon, and knowledge profile}
\begin{definition}[Predecessor trace]
\label{def:predecessor-trace}
The predecessor trace is the ordered family of prior cuts and their local readouts along a Memory path.
\end{definition}

\begin{definition}[Successor horizon]
\label{def:successor-horizon}
The successor horizon is the admissible family of next continuations generated from the current Memory State, the Markov Kernel, and the declared local constraints.
\end{definition}

\begin{choice}[Isotropic ignorance]
\label{choice:isotropic-ignorance}
The AFC isotropic-ignorance posture is used on the declared cut unless an explicit local refinement is declared.
\end{choice}

\begin{choice}[Anisotropic knowledge]
\label{choice:anisotropic-knowledge}
Anisotropic knowledge is a declared local refinement of distinctions or support on the declared cut. It narrows the admissible successor family on that cut. Dwell appears when a dwell condition is separately declared.
\end{choice}

\subsection{Inherited support and canonical Q4 chart}
\begin{choice}[Inherited support]
\label{choice:inherited-support}
The support law used below is the inherited AFC support law on the declared readout family.
\end{choice}

\begin{choice}[Canonical Q4 chart]
\label{choice:q4-chart}
At the chosen local resolution, the inherited induced support is read on a canonical Q4 chart.
\end{choice}

\begin{choice}[Axis-lift chart policy]
\label{choice:axis-lift}
A declared axis-lift map associates each active ternary step with a signed toggle of exactly one of four chart axes,
\[
\pi_{\mathrm{axis}}:(X_t,X_{t+1})\mapsto \Delta x_t\in\{\pm e_1,\pm e_2,\pm e_3,\pm e_4\},
\]
producing a Q4 chart representation of the inherited support. The chart is the declared row form used below.
\end{choice}

\begin{definition}[Q4 cell basis]
\label{def:q4-basis}
On the canonical chart the basic cell types are the Q4 vertices, Hamming-1 edges, square faces, cubes, and the 4-cell body. Traversal bookkeeping records per-axis signed counts on this same basis.
\end{definition}

\subsection{Loop-witness discipline and even-cycle typing}
\begin{definition}[Loop witness]
\label{def:loop-witness}
A loop witness is a closure witness on the inherited support.
\end{definition}

\begin{definition}[Two-cycle return, four-cycle enclosure, and repeating closure signature]
\label{def:cycle-typing}
A two-cycle return is the minimal return closure. A four-cycle enclosure is the minimal nondegenerate enclosure on the canonical chart. Repeating closure signatures are the declared even-cycle classifications built from these two witnesses.
\end{definition}

\begin{definition}[Level index and beat closure]
\label{def:level-beat}
For any declared closure witness $\kappa$ with active half-length count $L_{\pm}(\kappa)$, define the level index
\[
\lambda(\kappa):=\frac{L_{\pm}(\kappa)}{4}.
\]
The primitive beat closure is the minimal return witness with $L_{\pm}=2$. The first nondegenerate enclosure on the canonical Q4 chart has $L_{\pm}=4$.
\end{definition}

\subsection{Closure subclasses and minimal cage-band typing}
\begin{definition}[Closure predicate]
\label{def:closure-predicate}
A closure occurrence $\kappa$ fires on a segment $s$ when the active displacement closes on the declared readout,
\[
\sum_{t\in s}\tau_t=0,
\]
and the declared detector admissibility criteria hold on that same segment.
\end{definition}

\begin{definition}[Beat-knot and cage-knot]
\label{def:beat-cage-knot}
A knot is a closure occurrence treated as a localized curl on the wavelet stream. A \emph{beat-knot} is the minimal closure beat with $L_{\pm}=2$. A \emph{cage-knot} is a closure occurrence that additionally encloses a 2-cell under the canonical Q4 chart on the same declared detector policy.
\end{definition}

\begin{definition}[Cage as a two-knot object and minimal cage-band]
\label{def:cage-band}
A cage row is a two-knot object
\[
\mathrm{Cage}(M;\omega):=(\kappa_{\mathrm{in}},\kappa_{\mathrm{out}}),
\]
where $\kappa_{\mathrm{out}}$ is a cage-knot and $\kappa_{\mathrm{in}}$ is a knot nested inside $\kappa_{\mathrm{out}}$ under the canonical Q4 chart and the declared motif-to-region extractor. The minimal cage-band length signature is recorded as $(L_{\mathrm{in}}:L_{\mathrm{out}})=(4{:}6)$.
\end{definition}

\begin{definition}[Outer enclosure token and motif id]
\label{def:outer-enclosure-token}
The outer enclosure token $\eta_{\mathrm{out}}$ records the outer enclosure class on the declared window. A motif id $M$ is a stable label for a detected class of occurrences under declared detector settings and perturbation tolerances.
\end{definition}

\subsection{\texorpdfstring{$\alpha\leftrightarrow\omega$}{alpha-omega} bookkeeping, backlog, and crowning}
\begin{definition}[$\alpha$ tally and $\omega$ obligation]
\label{def:alpha-omega}
$\alpha$ denotes the local interior tally. $\omega$ denotes the closure or return obligation on the same cut.
\end{definition}

\begin{definition}[Backlog and reconciliation]
\label{def:backlog-reconciliation}
Backlog is the unresolved mismatch between the local tally and the closure obligation. Reconciliation is the declared balancing process that services that backlog.
\end{definition}

\begin{definition}[Crowning]
\label{def:crowning}
Crowning is the directional closure distinction that first stabilizes a local asymmetry under the declared crown predicate.
\end{definition}

\begin{definition}[Drift]
\label{def:drift}
Drift is the phase or transport bias that persists across the declared cut family on a declared readout class.
\end{definition}

\subsection{AO field, wavelet, shell/cadence interface, and saturation}
\begin{definition}[AO field]
\label{def:ao-field}
The AO field is the ordered bundle of local AOD quantities attached to a cut and its unspooling.
\end{definition}

\begin{definition}[Retained AO-field specification]
\label{def:retained-field-specification}
For an AO field \(\mathcal F\) on a declared window \(\omega\), a retained field specification is a declared tuple
\[
\mathfrak S
=
(\mathcal F,\mathcal F_K,\omega,\ell,\mathcal C_{\mathrm{omit}},\mathcal F_{K+1},\pi),
\]
where \(\mathcal F_K\subset\mathcal F\) is the retained extraction, \(\ell\) is the working cut level, \(\mathcal C_{\mathrm{omit}}\) records omitted context, \(\mathcal F_{K+1}\) is the first refinement when declared, and \(\pi\) is the local extraction/readout policy. This is bookkeeping on the declared AO field under the imported AFC basis.
\end{definition}

\begin{definition}[Refinement residual]
\label{def:refinement-residual}
For a retained extraction \(\mathcal F_K\subset\mathcal F\) and first refinement \(\mathcal F_{K+1}\), the refinement residual for a row functional \(Q\) on \(\omega\) is
\[
\widehat\epsilon_{K\to K+1}(\omega)
=
\left|
Q(\mathcal F_{K+1},\omega)-Q(\mathcal F_K,\omega)
\right|.
\]
A qualitative residual is a placeholder until the refinement row is computed.
\end{definition}

\subsection{Integer boundary-path row layer}


\begin{definition}[Exact row form]
\label{def:exact-row-form}
An AOD row records declared cut data as exact count, sign, residue, rational-pair, cycle-pair, and incidence data on a declared window. Its entries may include integer counts, signed differences, ternary signs, modular residues, support classes, phase-cycle pairs, path-count pairs, orientation signs, boundary incidences, and topology signatures.

The row is read at the declared working level \(\ell\) and window \(\omega\). A full phase closure is one cycle. Directional state is recorded by signed residue, orientation, and incidence data attached to the retained boundary relation.
\end{definition}

\begin{definition}[Cycle phase row]
\label{def:cycle-phase-row}
A retained boundary link may carry cycle data
\[
\Phi_e=(c_e,k_e,N_e),
\qquad
c_e\in\mathbb Z,
\quad
k_e\in\mathbb Z/N_e\mathbb Z,
\quad
N_e\in\mathbb Z_{>0}.
\]
Here \(c_e\) is the retained completed-cycle count when declared, \(k_e\) is the residual phase count, and \(N_e\) is the cycle modulus.

For adjacent endpoint residues define
\[
\Delta k_e:=\operatorname{wrap}_{N_e}(k_{e,+}-k_{e,-}),
\qquad
r_e:=\operatorname{cent}_{N_e}(\Delta k_e),
\]
and the direction sign
\[
d_e:=\operatorname{sgn}(r_e)\in\{-1,0,+1\}.
\]
The centered representative \(r_e\) is the phase-step count used by the retained transport row.
\end{definition}

\begin{definition}[Relative temporal-duration transport row]
\label{def:exact-path-count-ratio}
A retained transport row records relative temporal-duration change as the integer pair
\[
v_e^{\mathrm{exact}}=(\Delta s_e,\Delta n_e),
\qquad
\Delta s_e,\Delta n_e\in\mathbb Z,
\qquad
\Delta n_e>0,
\]
where \(\Delta s_e\) is a path-step, shell-step, boundary-hop, or topology-step count, and \(\Delta n_e\) is a duration count in \(\mathrm{bip}_{\ell}\). Direction is carried by the incidence sign on the same retained relation.
\end{definition}

\begin{definition}[Effective-charge exact pair]
\label{def:qeff-exact-pair}
For a retained boundary link or face \(e\), let
\[
Q_e:=J_{\alpha}(e;\omega)-J_{\Omega}(e;\omega)\in\mathbb Z,
\]
\[
S_e:=J_{\alpha}(e;\omega)+J_{\Omega}(e;\omega)+J_0(e;\omega)\in\mathbb Z_{>0}.
\]
The exact effective-charge row is
\[
q_{\mathrm{eff},e}^{\mathrm{exact}}=(Q_e,S_e).
\]
The pair records chirality-weighted boundary polarization counts on the declared cut.
\end{definition}

\begin{definition}[Rational-pair arithmetic]
\label{def:rational-pair-arithmetic}
For integer pairs with positive denominators, multiplication and addition are
\[
(a,b)\cdot(c,d):=(ac,bd),
\qquad
(a,b)+(c,d):=(ad+bc,bd),
\qquad b,d>0.
\]
The notation \(\operatorname{Rat}[a,b]\) denotes the reduced integer pair with positive denominator. If \(a=0\), its canonical form is \((0,1)\).
\end{definition}

\begin{definition}[Integer attenuation pairs]
\label{def:integer-attenuation-pairs}
For a retained boundary link \(e\), the support attenuation pair is
\[
A_{B,e}^{\mathrm{exact}}
=(A_{B,e}^{+},A_{B,e}^{-})
=(B_0^{\widehat\Lambda_e},(B_e^*)^{\widehat\Lambda_e}).
\]
When \(B_0=1\), this is \((1,(B_e^*)^{\widehat\Lambda_e})\). The loss attenuation pair is
\[
A_{R,e}^{\mathrm{exact}}
=(A_{R,e}^{+},A_{R,e}^{-})
=(R_{0,e},R_{0,e}+R_e),
\qquad
R_{0,e}\in\mathbb Z_{>0}.
\]
The total attenuation gate is
\[
A_e^{\mathrm{exact}}=A_{B,e}^{\mathrm{exact}}\cdot A_{R,e}^{\mathrm{exact}}.
\]
Here \(B^*\), \(\Lambda_b\), and \(\widehat\Lambda\) name distinct quantities: \(B^*\) records the dominant support count, \(\Lambda_b\) records support mass on the dominant class, and \(\widehat\Lambda\) records the attenuation exponent used by this boundary row.
\end{definition}

\begin{definition}[Torsion and AOD momentum proxy pairs]
\label{def:torsion-momentum-proxy-pairs}
The torsion gain pair on a retained boundary link is
\[
G_{T,e}^{\mathrm{exact}}=(T_{0,e}+T_e,T_{0,e}),
\qquad
T_{0,e}>0.
\]
The positive numerator branch uses \(T_{0,e}+T_e>0\); signed numerator pairs are declared when the row requires them. A momentum-like row is recorded as a named AOD proxy,
\[
p_e^{\mathrm{AOD}}=(P_e,P_{0,e}),
\qquad
P_e,P_{0,e}\in\mathbb Z,
\qquad
P_{0,e}>0.
\]
The map that produces \(p_e^{\mathrm{AOD}}\) is declared before the proxy is used.
\end{definition}

\begin{definition}[Retained boundary-link row]
\label{def:retained-boundary-link-row}
A retained boundary link on level \(\ell\) and window \(\omega\) may be recorded as
\[
D_e^{(\ell,\omega)}=
\left(
 c_e,k_e,N_e,
 \chi_e,
 Q_e,S_e,
 R_e,R_{0,e},
 T_e,T_{0,e},
 B_e^*,\widehat\Lambda_e,
 P_e,P_{0,e},
 \Sigma_e
\right).
\]
The entries record, in order, cycle data, chirality, effective-charge counts, loss counts, torsion pair, support-count witness, attenuation exponent, AOD momentum proxy, and topology signature. This is a retained boundary-link row on a declared AO field.
\end{definition}

\begin{definition}[Topology signature]
\label{def:topology-signature}
The topology signature of a retained boundary relation is the discrete object
\[
\Sigma_e=(\tau_e,\sigma_e,\mathbf a_e,\partial_e,\mathcal I_e),
\]
where \(\tau_e\) is a cell or face type index, \(\sigma_e\in\{-1,0,+1\}\) is an orientation sign, \(\mathbf a_e\in\mathbb Z^m\) is a fractal or scale address, \(\partial_e\) is a boundary incidence list, and \(\mathcal I_e\) is an integer incidence matrix or tensor. Incidence signs are read as
\[
\epsilon_{ej}=\epsilon(\Sigma_e,\Sigma_j)\in\{-1,0,+1\}.
\]
The topology signature supplies the incidence signs used by retained boundary relations.
\end{definition}

\begin{definition}[Exact transport flux]
\label{def:exact-transport-flux}
For a retained boundary link \(e\), set
\[
\Delta\Phi_e=(r_e,N_e),
\qquad
q_e=(Q_e,S_e),
\]
with attenuation pairs \(A_{B,e}\), \(A_{R,e}\), and torsion pair \(G_{T,e}=(T_{0,e}+T_e,T_{0,e})\). The retained transport flux is
\[
u_e^{(\mathrm{tr})}
=
\operatorname{Rat}
\left[
 r_e Q_e A_{B,e}^{+}A_{R,e}^{+}(T_{0,e}+T_e),
 N_e S_e A_{B,e}^{-}A_{R,e}^{-}T_{0,e}
\right].
\]
Its sign is \(\operatorname{sgn}(u_e^{(\mathrm{tr})})=\operatorname{sgn}(r_eQ_e)\). Chirality is applied in the transverse incidence readout.
\end{definition}

\begin{definition}[Windowed current-flux rate]
\label{def:windowed-current-flux-rate}
Let \(\epsilon_e^{\parallel}\) be the retained-path current-flux direction. Let \(\epsilon_{ej}^{\perp}\) be the transverse probe/readout incidence sign. These signs occupy separate longitudinal and transverse incidence slots. The accumulated retained transport is
\[
U_K^{(\ell)}(\omega)
=
\sum_{e\in E_K(\omega)}\epsilon_e^{\parallel}u_e^{(\mathrm{tr})}.
\]
For \(|\omega|_{\mathrm{bip}_{\ell}}\in\mathbb Z_{>0}\), the windowed current-flux rate is
\[
\mathcal J_K^{(\ell)}(\omega)
=
\operatorname{Rat}
\left[
U_{K,\mathrm{num}}^{(\ell)}(\omega),
U_{K,\mathrm{den}}^{(\ell)}(\omega)|\omega|_{\mathrm{bip}_{\ell}}
\right].
\]
The symbol \(\mathcal J_K\) names the windowed retained current-flux rate and is separate from the chirality-count rows \(J_{\alpha}\), \(J_{\Omega}\), and \(J_0\).
\end{definition}

\begin{definition}[Carrier-family partition]
\label{def:carrier-family-partition}
For the first retained transport implementation, carrier families are declared as a disjoint partition
\[
E_K(\omega)=E_{K,\mathrm{duon}}(\omega)\cup E_{K,\mathrm{tetron}}(\omega)\cup E_{K,\mathrm{other}}(\omega),
\]
\[
E_{K,a}\cap E_{K,b}=\varnothing\qquad(a\neq b).
\]
Thus \(U_K\) and \(\mathcal J_K\) split exactly into duon, tetron, and other carrier contributions.
\end{definition}

\begin{definition}[Transverse boundary witness]
\label{def:transverse-boundary-witness}
Let \(\epsilon_{ej}^{\perp}\in\{-1,0,+1\}\) be the transverse incidence sign from retained link \(e\) to probe side \(j\). The exact transverse witness is
\[
w_{\perp,j}^{\mathrm{exact}}
=
\sum_e \epsilon_{ej}^{\perp}\chi_e u_e^{(\mathrm{tr})}.
\]
For the two-sided wire witness,
\[
\epsilon_{e,+}^{\perp}=+1,
\qquad
\epsilon_{e,-}^{\perp}=-1,
\]
so
\[
w_{\perp}^{+}=\sum_e\chi_eu_e^{(\mathrm{tr})},
\qquad
w_{\perp}^{-}=-\sum_e\chi_eu_e^{(\mathrm{tr})}.
\]
The transverse boundary witness is denoted \(w_{\perp}\), and the support witness \(B^*\), current-flux rate \(\mathcal J_K\), and transverse witness \(w_{\perp}\) are pairwise distinct:
\[
\boxed{B^*\neq\mathcal J_K\neq w_{\perp}}.
\]
\end{definition}

\begin{definition}[Integer retained-boundary diagram]
\label{def:integer-retained-boundary-diagram}
A retained-boundary or retained-path diagram value is
\[
\mathcal A_{\mathbb Z}(\mathcal G)
=(W_{\mathcal G},K_{\mathcal G},N_{\mathcal G},\chi_{\mathcal G},\Sigma_{\mathcal G}),
\]
where \(W_{\mathcal G}\) is a rational weight pair, \(K_{\mathcal G}\) is an accumulated phase residue, \(N_{\mathcal G}\) is the accumulated modulus, \(\chi_{\mathcal G}\) is an accumulated chirality sign, and \(\Sigma_{\mathcal G}\) is the accumulated topology signature. For a path \(\gamma\),
\[
K_{\gamma}=\sum_{e\in\gamma}k_e\pmod{N_{\gamma}}.
\]
The retained-path accumulation is
\[
\mathcal Z_{\partial\mathcal O}^{(\ell,\omega)}
=
\bigoplus_{\gamma\in\Gamma_{\partial\mathcal O}(\omega)}
(W_{\gamma},K_{\gamma},N_{\gamma},\chi_{\gamma},\Sigma_{\gamma}).
\]
The retained-path accumulation is taken over the admissible retained path family.
\end{definition}

\begin{definition}[Integer collision vertex]
\label{def:integer-collision-vertex}
For incoming boundary signs \(a,b\in\{-1,0,+1\}\), the balanced ternary collision update is \(c=a\oplus_3b\). A retained collision vertex is
\[
\mathcal K_{\mathrm{coll}}^{\mathbb Z}:
D_a\times D_b\times\Sigma_{ab}
\longrightarrow
D'\oplus\Pi_{\mathrm{export}}\oplus\mathcal R_{\mathrm{cut}}.
\]
The update is controlled by integer predicates
\[
a\oplus_3 b,
\qquad
Q_a+Q_b,
\qquad
K_a+K_b\pmod N,
\qquad
\chi_a\chi_b,
\qquad
\epsilon(\Sigma_a,\Sigma_b).
\]
The outcome class is one of balanced, push, pull, reclosure, export, or rupture/refinement.
\end{definition}

\begin{definition}[Wavelet]
\label{def:wavelet}
A wavelet is a finite closure-carrying Memory-path segment equipped with ordered cutwise field quantities.
\end{definition}

\begin{definition}[Shell-band extraction interface]
\label{def:shell-cadence-interface}
A shell-band extractor packages the declared closure/return/cycle data into an ordered family of shell-band rows on the same declared window.
\end{definition}

\begin{definition}[Cadence extractor]
\label{def:cadence-extractor}
A cadence extractor acts on the declared shell bands to produce completed-cycle counts, completion periods, phase indices, cycle-ratio phase rows, and drift witnesses on the same declared window.
\end{definition}

\begin{definition}[Saturation]
\label{def:saturation}
Saturation is the occupancy of declared traversal lanes or closure seats on the local readout.
\end{definition}

\subsection{Push/pull, effective charge, and export/radiation typing}
\begin{definition}[Push/pull boundary current]
\label{def:push-pull-current}
Push/pull is the chirality-resolved boundary-current bookkeeping read from the same oriented boundary histories used for lag and torsion. The sign convention must be declared once for the declared readout family.
\end{definition}

\begin{definition}[Effective-charge proxy]
\label{def:qeff}
With chirality-resolved boundary totals
\[
J_\alpha(M;\omega):=\sum_{t\in\omega}\sum_{(i,j)\in\partial R(M;\omega)} \phi_{ij}(t)\,\mathbf 1\{\chi(t)=\alpha\},
\]
\[
J_\Omega(M;\omega):=\sum_{t\in\omega}\sum_{(i,j)\in\partial R(M;\omega)} \phi_{ij}(t)\,\mathbf 1\{\chi(t)=\Omega\},
\]
and neutral or uncommitted count \(J_0(M;\omega)\) when the row declares it, define
\[
Q(M;\omega):=J_\alpha(M;\omega)-J_\Omega(M;\omega)\in\mathbb Z,
\]
\[
S(M;\omega):=J_\alpha(M;\omega)+J_\Omega(M;\omega)+J_0(M;\omega)\in\mathbb Z_{>0}.
\]
The exact effective-charge proxy is
\[
q_{\mathrm{eff}}^{\mathrm{exact}}(M;\omega):=(Q(M;\omega),S(M;\omega)).
\]
A declared scalar readout is obtained by the declared scalar readout map \(\mathcal N(Q,S)\).
\end{definition}

\begin{definition}[$\Omega$-cycle lag]
\label{def:phase-lag}
The chirality-resolved cycle-lag statistic is the declared difference between paired $\Omega$ and $\alpha$ cycle-ratio histories on the declared readout family. Its pairing policy must be explicit.
\end{definition}

\begin{definition}[Radiating/export window and emitted packet witness]
\label{def:radiating-window}
A window is radiating/exporting for motif $M$ when the declared shell predicate records nonzero net export and/or elevated boundary mutual information on the same readout family. The emitted/exported packet witness is the wavelet packet crossing the declared cut on that window.
\end{definition}

\subsection{\texorpdfstring{$B^*$}{B*} extraction, histogram support, and Jensen discipline}
\begin{definition}[Return-duration extraction]
\label{def:b-extraction}
For each declared return or closure occurrence $\kappa$, the return-duration extraction policy emits a burden sample
\[
b(\kappa):=e_\alpha(\kappa)+e_\Omega(\kappa)+f(\kappa),\qquad f(\kappa)\in\mathbb Z_{\ge 1},
\]
where $e_\alpha,e_\Omega$ are the one-way outbound/return durations and $f$ is the measured bridge/reflection count under the declared policy.
\end{definition}

\begin{definition}[Histogram, support mass, and dominant residual-duration class]
\label{def:bstar-histogram}
Let $\mathrm{Occ}(M;\omega)$ denote the declared set of motif occurrences. With declared weights $w(\kappa)\ge 0$ (TTL weighting allowed but not required), define the residual-duration histogram
\[
H_B(b;M,\omega):=\sum_{\kappa\in\mathrm{Occ}(M;\omega)} w(\kappa)\,\mathbf 1\{b(\kappa)=b\},
\]
the dominant residual-duration class
\[
B^*(M;\omega):=\arg\max_b H_B(b;M,\omega),
\]
and the support mass on that class
\[
\Lambda_b(M;\omega):=H_B(B^*(M;\omega);M,\omega).
\]
The estimator functional must be declared whenever ties or trimming rules are present.

The value \(B^*(M;\omega)\) is a support-count quantity recorded in \(\mathrm{bip}_0\) at the declared working cut level. Its reciprocal \(1/B^*(M;\omega)\), when used as a rate quantity on the same declared row, is recorded in \(\mathrm{biz}_0\). The subscript \(0\) marks the declared reference level; see Appendix~\ref{app:level-relative-bip-biz}.
\end{definition}

\begin{remarkx}
\label{rem:bstar-class}
The $\alpha/\Omega$ branch mismatch generates the residual-duration burden on the declared window. $B^*(M;\omega)$ is the dominant class of that burden; $\Lambda_b(M;\omega)$ is the support mass on that class.
\end{remarkx}

\begin{choice}[Minimal-pivot identity bookkeeping]
\label{choice:high-rung-pivot}
A declared minimal-pivot identity may be used to expose the integer accounting spine. Under that idealization one sets
\[
b=2e+f,\qquad f\in\mathbb Z_{\ge1}.
\]
\end{choice}

\begin{remarkx}[Jensen discipline]
\label{rem:jensen}
Histogram and loss aggregation must preserve the order of operations. If $T_{\mathrm{eval}}=n_1+n_2+n_3$ and
\[
\bar m:=\frac{n_1+2n_2+3n_3}{T_{\mathrm{eval}}},
\]
then the exact loss-average on the declared window is
\[
R_{\mathrm{loss}}(\omega)=\frac{n_1+n_2\log_3\!\left(\frac{3}{2}\right)}{T_{\mathrm{eval}}},
\]
whereas the post-averaged surrogate would be
\[
\log_3\!\left(\frac{3}{\bar m}\right).
\]
Their order-sensitive difference is recorded by
\[
J_{\mathrm{gap}}(\omega):=
\frac{n_1+n_2\log_3\!\left(\frac{3}{2}\right)}{T_{\mathrm{eval}}}
-
\log_3\!\left(\frac{3}{\bar m}\right)\ge 0,
\]
The declared loss row aggregates on the sample set before later compression.
\end{remarkx}

\begin{choice}[Shell-band extractor]
\label{choice:shell-band-extractor}
A shell-band extractor packages the declared closure, return, and dressing stack into an ordered family of shell bands on the same window. The extractor policy id must be declared and stable across reruns.
\end{choice}

\begin{definition}[Shell signature]
\label{def:shell-signature}
On a declared motif-conditioned window $\omega$, the shell signature is the coarse fingerprint
\[
\mathrm{Sig}(M;\omega):=(\lambda\text{-band}, B^*\text{-class}, \rho\text{-winding class}, V\text{-class}).
\]
It is assembled from the level index, dominant residual-duration class, winding row, and support signature on the same declared shell-band extraction policy.
\end{definition}

\begin{definition}[Confinement predicate and proximity criterion]
\label{def:confinement-predicate}
A motif-conditioned window is declared confined when it is shell-qualified, carries the declared $B^*$-dominance on the same shell-signature class, and satisfies the declared proximity criterion for the relevant traverser/closure family together with the corresponding cage-band / outer-enclosure requirement on the same declared window. The transmutation trigger is the declared transition condition paired with that confinement predicate on the same declared window.
\end{definition}

\subsection{Traverser / closure naming rule, field aliases, and structural constraints}
\begin{definition}[Traverser / closure rule]
\label{def:traverser-closure-rule}
The suffix \texttt{-on} denotes a traverser class. The suffix \texttt{-ad} denotes a closure or seat class. Traverser precedes closure in composed names.
\end{definition}

\begin{definition}[Public AO field families and structural key]
\label{def:field-families}
The public AO field families used below are the two-seat family \emph{duon} and the four-seat family \emph{tetron}, together with their declared closure aliases \emph{duad} and \emph{tetrad} on the same family keys, declared structural subclasses, and O-prefixed chirality tags on the same field row. Here \emph{duad} is the closed/full duon alias and \emph{tetrad} is the locked/full tetron alias on the declared window. Structural examples such as \texttt{tetron:1:2:6}, \texttt{tetron:1:2:8}, \texttt{tetron:1:2:10}, \texttt{tetron:3:3:6}, and \texttt{tetrad:3:4:6} are public identifiers on declared windows.
\end{definition}

\begin{definition}[Structural constraints]
\label{def:structural-constraints}
Seat / traversal constraints, closure-lock constraints, and saturation policy are declared on the declared family key before any subclass persistence or transmutation claim is made.
\end{definition}

\begin{definition}[Ominion]
\label{def:ominion}
An ominion is a declared clustered effective-body state on a declared window, carrying effective row properties such as mass/rate, charge proxy, linear momentum, angular momentum, phase state, and loss/tension witnesses.
\end{definition}

\begin{definition}[n-soma system]
\label{def:n-soma}
An n-soma is a finite family of ominions together with a shared relational record on the declared window.
\end{definition}

\subsection{Observer families, coupling, chirality, polyrotor phases, torsional tension, and cadence}
\begin{definition}[Observer family and coupling]
\label{def:observer-coupling}
An observer family is a declared family of Memory-State traces. Coupling is a shared narrowing relation or path convergence on common successor support.
\end{definition}

\begin{definition}[Chirality]
\label{def:chirality}
Chirality is the persistent handedness of asymmetry under unspooling.
\end{definition}

\begin{definition}[Polyrotor phase family]
\label{def:polyrotor}
A polyrotor phase family is a finite indexed family of rotor-completion readouts on the same declared cut family.
\end{definition}

\begin{definition}[Rotor-completion period, phase index, cycle-ratio phase, and completed-cycle count]
\label{def:rotor-completion}
On the declared shell/window, let $P_\rho(M;\omega)\in \mathbb N_{>0}$ be the rotor-completion period, let $\rho_t(M;\omega)\in\{0,1,\dots,P_\rho(M;\omega)-1\}$ be the phase index, let
\[
\varphi_t(M;\omega):=\frac{\rho_t(M;\omega)}{P_\rho(M;\omega)}\in[0,1)\cap\mathbb Q
\]
be the cycle-ratio phase, and let $N_\rho(M;\omega)\in\mathbb Z_{\ge 0}$ be the cumulative completed-cycle count. A complete cycle increments $N_\rho$ by one.
\end{definition}

\begin{definition}[Torsional tension and torsional loss]
\label{def:torsional}
On a declared coupled polyrotor phase family, torsional tension is the accumulated unresolved phase disagreement carried on that family. Torsional loss is the admissibility lost because that disagreement remains unresolved.
\end{definition}

\begin{definition}[Recursive phase burden]
\label{def:recursive-phase-debt-service}
On a declared coupled polyrotor phase family, recursive phase debt service is the internal servicing of unresolved phase disagreement under the recursive update. On the declared readout family it is recorded through torsional tension, torsional loss, the torsional residual witness, and any forced dwell where present.
\end{definition}

\begin{definition}[Torsional residual witness]
\label{def:torsional-residual-witness}
The torsional residual witness is the noncancelling remainder of the chirality-resolved boundary-current histories on the declared readout family. Its explicit operator realization is recorded later in the operator section.
\end{definition}

\begin{definition}[Cadence]
\label{def:cadence-bstar}
Cadence is the declared local recurrence-rate quantity on the declared readout class. The dominant residual-duration class $B^*$ is defined earlier by the declared return-duration histogram and is reused here only by reference.
\end{definition}


\subsection{Field persistence, decay, and recruitment}
\begin{definition}[TTL hazard and right-censoring]
\label{def:ttl-hazard}
For a declared motif id $M$, detector policy, and rupture predicate $H$, the lifetime observable is the birth--death duration
\[
\mathrm{TTL}_{\mathrm{bip}_0}(M;\omega):=t_{\mathrm{death}}-t_{\mathrm{birth}}\in\mathbb N,
\]
with right-censoring declared when the declared window ends before rupture. TTL weighting or filtering is attached by explicit policy id.
\end{definition}

\begin{definition}[Decay, by-product census, and rupture predicate]
\label{def:decay-census}
Decay is the typed end of persistence under the declared rupture predicate. Product and by-product census are the typed post-death motif counts on declared post-death windows. These are row objects on the same underlying AO field.
\end{definition}

\begin{definition}[Freezing, rebalance dwell, transmutation, and dynamic recruitment]
\label{def:freeze-recruit}
A declared window is \emph{freezing} when sustained admissibility restriction and/or sustained CK mismatch indicates that the current representation is saturating the imported AFC Markov-kernel capacity $C(P)$ on that representation. Rebalance events are explicit operator-logged events with nonzero dwell cost. Transmutation is a change of motif signature class across windows with explicit dwell cost. Dynamic volume recruitment is the growth of effective closure range or outer enclosure used to preserve admissibility and shell qualification under continued load.
\end{definition}

\begin{definition}[Relative temporal-duration change and rest-frame window]
\label{def:rest-frame-window}
Using the declared axis-traversal rows, define the relative temporal-duration change
\[
v(\omega):=\frac{\|\Delta R(\omega)\|_2}{T_{\mathrm{act}}(\omega)}\in[0,1].
\]
The traversal direction is carried by the declared orientation/incidence data on the same row. A window is rest-frame qualified when it is shell-qualified and $v(\omega)\le \varepsilon_v$ under the declared tolerance.
\end{definition}
